% ============================================================================
% V2P Video Agent — Working Paper Draft
% ============================================================================
\documentclass[10pt,twocolumn]{article}

% --- packages ---------------------------------------------------------------
\usepackage[margin=1in]{geometry}
\usepackage{graphicx}
\usepackage{amsmath,amssymb}
\usepackage{booktabs}
\usepackage{algorithm}
\usepackage{algpseudocode}
\usepackage{hyperref}
\usepackage{xcolor}
\usepackage{enumitem}
\usepackage{caption}
\usepackage{subcaption}
% TikZ removed: figures are included from pre-rendered PDFs.
% Convert SVGs in docs/images/ to PDF before compiling:
%   for f in ../images/*.svg; do
%     inkscape "$f" --export-type=pdf --export-filename="$(basename "$f" .svg).pdf"
%   done
% Or install the 'svg' package and compile with --shell-escape.
\graphicspath{{../images/}}

\hypersetup{
  colorlinks=true,
  linkcolor=blue!60!black,
  citecolor=green!50!black,
  urlcolor=blue!70!black,
}

% convenience
\newcommand{\ie}{\textit{i.e.}}
\newcommand{\eg}{\textit{e.g.}}
\newcommand{\etal}{\textit{et al.}}
\newcommand{\sys}{\textsc{V2P Video Agent}}

% ============================================================================
\title{\sys: An Agentic Framework for Scalable\\
       Robot Demonstration Video Analysis}

\author{
  V2P Team \\
  NVIDIA \\
  \texttt{v2p@nvidia.com}
}

\date{\today}

% ============================================================================
\begin{document}
\maketitle

% ============================================================================
% ABSTRACT
% ============================================================================
\begin{abstract}
Learning robot manipulation policies from human demonstration videos
requires large-scale, accurately annotated datasets of atomic actions,
the objects involved, and the temporal boundaries of each interaction.
Producing such annotations manually is prohibitively expensive.
We present \sys{}, an end-to-end agentic framework that
(i)~segments raw demonstration videos into atomic manipulation clips
  using a Vision--Language Model (VLM) with a self-correcting
  verification--refinement loop,
(ii)~constructs a temporally grounded entity--relationship graph
  capturing the objects, people, and locations that appear across the
  video,
(iii)~indexes frame-level visual embeddings for semantic retrieval, and
(iv)~provides an agentic retrieval system that, given a high-level
  task description (\eg, ``make me a coffee''), decomposes it into
  sub-tasks and retrieves the most relevant training clips using a
  strict-to-relaxed search hierarchy.
The entire system is orchestrated with LangGraph and scales to
hundreds of videos via multi-GPU parallel processing with
duration-aware Longest-Processing-Time (LPT) sharding.
We evaluate the ingestion pipeline on EPIC-KITCHENS-100 using a
two-tier metric suite covering temporal segmentation quality
(mAP@tIoU, boundary F1) and annotation accuracy (top-$k$ class
matching, semantic similarity, BERTScore).
\end{abstract}

% ============================================================================
% 1  INTRODUCTION
% ============================================================================
\section{Introduction}\label{sec:intro}

Imitation learning and policy distillation from human demonstrations
are among the most practical paths toward general-purpose robot
manipulation~\cite{chi2023diffusion,zhao2023aloha,brohan2023rt2}.
A critical bottleneck is the \emph{data preparation} step:
raw demonstration videos---often hours long---must be temporally
segmented into short, atomic action clips, each annotated with the
manipulated object, the action performed, and precise temporal
boundaries.
Manual annotation does not scale; a single hour of video can take
an expert annotator several hours to label.

Recent advances in Vision--Language Models (VLMs) such as
Cosmos-Reason~\cite{nvidia2025cosmos} and
Qwen-VL~\cite{bai2023qwen} make it feasible to automate this
process.  However, naively prompting a VLM once over a long video
yields noisy, inconsistent segmentations.
Reliable automation requires \emph{closed-loop verification}
(a ``critic'' pass), \emph{structured knowledge extraction}
(entity graphs), and \emph{scalable orchestration} across
many videos and GPUs.

In this paper we describe \sys{}, an agentic framework that
addresses these challenges.  Our contributions are:

\begin{enumerate}[leftmargin=*,itemsep=2pt]
  \item \textbf{Self-correcting ingestion pipeline.}
        A LangGraph-based pipeline that segments videos via
        chunk-based VLM processing, verifies each clip with a
        VLM critic, and iteratively refines invalid segments---all
        with configurable deduplication (heuristic or LLM-based).

  \item \textbf{Entity graph construction.}
        An LLM-driven extraction stage that builds a temporally
        grounded entity--relationship graph (persons, objects,
        locations and their interactions), linked across overlapping
        chunks via time-gap-based entity resolution.

  \item \textbf{Agentic retrieval with strict-to-relaxed search.}
        A retrieval agent that decomposes high-level task queries
        into sub-tasks, executes them in parallel via the LangGraph
        Send API, and applies a four-level search relaxation
        hierarchy over the entity graph and frame embeddings.

  \item \textbf{Scalable multi-GPU architecture.}
        Duration-aware LPT sharding, chunked video processing,
        resume-capable batch ingestion, and concurrent SQLite writes
        (WAL mode) enable processing of large video corpora on
        multi-GPU nodes.
\end{enumerate}

% ============================================================================
% 2  RELATED WORK
% ============================================================================
\section{Related Work}\label{sec:related}

\paragraph{Temporal action segmentation.}
Temporal Action Localization (TAL) methods such as
ActionFormer~\cite{zhang2022actionformer} and
TriDet~\cite{shi2023tridet} detect action boundaries in video.
These approaches typically require pre-extracted features and
class-specific training.
By contrast, our system uses a VLM in a zero-shot setting,
producing free-text action and object labels without any
task-specific training.

\paragraph{Scene and entity graphs from video.}
Scene graph generation from video has been explored for visual
relationship detection~\cite{ji2020action,cong2021spatial}.
EGAgent~\cite{xu2025egagent} introduced LLM-driven entity graph
construction from egocentric video captions, providing structured
representations of entities and their interactions over time.
Our entity extraction module adopts the prompt structure of
EGAgent while extending it with cross-chunk entity linking and
embedding-based visual grounding.

\paragraph{VLM/LLM-based video understanding.}
Large vision--language models have been applied to video captioning
\cite{li2023videochat,lin2023video}, temporal grounding
\cite{lei2021detecting}, and question answering
\cite{xu2017video}.
Our work leverages VLMs for both \emph{generation}
(segmentation, entity extraction) and \emph{verification}
(critic pass), forming a closed-loop agentic pipeline rather than
a single-pass inference.

\paragraph{Agentic workflows.}
Tool-augmented LLM agents~\cite{schick2023toolformer,yao2023react}
and multi-agent frameworks~\cite{wu2023autogen} have gained
traction for complex reasoning tasks.
LangGraph~\cite{langgraph2024} provides a graph-based orchestration
layer with support for conditional routing, cycles, and parallel
execution.  We build both the ingestion and retrieval pipelines
as LangGraph state machines, enabling the verification--refinement
loop and parallel sub-task execution.

\paragraph{Learning from demonstration videos.}
Video-conditioned policy learning has been explored through
approaches such as
RT-2~\cite{brohan2023rt2},
Diffusion Policy~\cite{chi2023diffusion}, and
ALOHA~\cite{zhao2023aloha}.
These methods assume curated clip datasets.
Our framework sits upstream, automating the curation of such
datasets from raw, unstructured demonstration recordings.


% ============================================================================
% 3  SYSTEM ARCHITECTURE
% ============================================================================
\section{System Architecture}\label{sec:arch}

The system consists of two main agentic
pipelines---\emph{Ingestion} (Figure~\ref{fig:ingestion-pipeline})
and \emph{Retrieval} (Figure~\ref{fig:retrieval-agent})---connected
through a dual-database layer (entity graph DB and frame vector DB),
with a scalable multi-GPU execution layer
(Figure~\ref{fig:scaling}).

\begin{figure*}[t]
  \centering
  \includegraphics[width=0.75\textwidth]{ingestion_langgraph.pdf}
  \caption{%
    \textbf{Ingestion pipeline (LangGraph state graph).}
    The pipeline segments a video via chunked VLM processing,
    optionally verifies each clip with a VLM critic (orange nodes),
    iteratively refines invalid segments, then constructs an entity
    graph (green nodes) with frame embeddings and writes to the
    graph and vector databases.
    Conditional edges enable toggling verification, refinement,
    entity graph construction, and reporting independently.
  }
  \label{fig:ingestion-pipeline}
\end{figure*}


% ---------------------------------------------------------------------------
\subsection{Ingestion Pipeline}\label{sec:ingestion}

The ingestion pipeline transforms a raw demonstration video into
(a)~a set of verified, atomic manipulation clips and
(b)~a structured entity--relationship graph stored in a SQLite
database.  The pipeline is implemented as a LangGraph
\texttt{StateGraph} with the following stages.

\subsubsection{Chunk-Based Video Segmentation}

Long videos are divided into overlapping temporal chunks of
configurable size $C$ seconds with overlap $O$ seconds
(defaults: $C{=}15$, $O{=}1.5$).
Each chunk is extracted to a temporary file using FFmpeg and
independently processed by a VLM (\eg, Cosmos-Reason2-8B or
Qwen3-VL-8B).

For each chunk, the VLM receives a structured prompt instructing it
to identify distinct manipulation actions and return a JSON array of
clips, each with \texttt{start\_time}, \texttt{end\_time},
\texttt{object}, \texttt{action}, and \texttt{description}.
Timestamps returned by the VLM are relative to the chunk start and
are converted to absolute video time during parsing.

Because adjacent chunks overlap, clips may be duplicated.
A deduplication stage merges temporally overlapping clips using
either a heuristic rule (keep the longer clip) or an LLM-based
decision that compares the two annotations and merges only when
both object and action match.

\subsubsection{VLM Critic Verification Loop}

Each candidate clip is extracted to a short \texttt{.mp4} file and
passed to a VLM ``critic'' that evaluates:
(i)~whether the clip contains a single coherent action,
(ii)~boundary quality (start and end),
(iii)~object identification accuracy, and
(iv)~action description correctness.
The critic returns a structured JSON verdict with a confidence
score and specific violation descriptions.

Invalid clips are routed to a \emph{refinement node} that
re-annotates them.  This verify--refine loop runs for up to
$K$ iterations (default $K{=}3$), with each iteration only
re-verifying previously invalid clips.
Algorithm~\ref{alg:verify-refine} summarises the procedure.

\begin{algorithm}[t]
\caption{Verify--Refine Loop}\label{alg:verify-refine}
\begin{algorithmic}[1]
\Require Clips $\mathcal{C}$, max iterations $K$
\State $\mathcal{V} \gets \textsc{VerifyAll}(\mathcal{C})$
\For{$k = 1, \dots, K$}
  \State $\mathcal{C}_{\text{inv}} \gets \{c \in \mathcal{C} \mid \mathcal{V}[c].\text{is\_valid} = \texttt{false}\}$
  \If{$\mathcal{C}_{\text{inv}} = \emptyset$}
    \State \textbf{break}
  \EndIf
  \State $\mathcal{C}_{\text{inv}} \gets \textsc{Refine}(\mathcal{C}_{\text{inv}})$
    \Comment{Re-annotate via VLM}
  \State Update $\mathcal{C}$ with refined clips
  \State $\mathcal{V}_{\text{new}} \gets \textsc{Verify}(\mathcal{C}_{\text{inv}})$
  \State Merge $\mathcal{V}_{\text{new}}$ into $\mathcal{V}$
\EndFor
\State \Return $\mathcal{C}$, $\mathcal{V}$
\end{algorithmic}
\end{algorithm}


% ---------------------------------------------------------------------------
\subsection{Entity Graph Construction}\label{sec:entity-graph}

After segmentation, the pipeline constructs a temporally grounded
entity--relationship graph.

\subsubsection{Entity and Relationship Extraction}

For each verified clip, the clip's textual description is passed to
an LLM (which may be the same VLM or a separate text-only model).
Following the prompt structure of EGAgent~\cite{xu2025egagent},
the model extracts:
\begin{itemize}[leftmargin=*,itemsep=1pt]
  \item \textbf{Entities} with fields: \texttt{id} (normalised
    name), \texttt{type} $\in \{\text{person}, \text{object},
    \text{location}\}$, and optional properties.
  \item \textbf{Relationships} with fields: \texttt{source},
    \texttt{target}, and \texttt{type} $\in \{\text{picks-up},
    \text{places}, \text{grasps}, \text{interacts-with},
    \ldots\}$.
\end{itemize}
Each entity and relationship inherits the temporal extent
$[t_{\text{start}}, t_{\text{end}}]$ of the clip it was extracted
from.

\subsubsection{Cross-Chunk Entity Linking}

Because the same physical object may appear in multiple clips,
entities with identical normalised IDs are grouped and merged based
on temporal proximity.  Specifically, two entity mentions $e_i$ and
$e_j$ with the same \texttt{entity\_id} are merged if
\begin{equation}
  e_j.\text{first\_seen} - \max_{e \in \text{cluster}}
    e.\text{last\_seen} \leq \tau,
\end{equation}
where $\tau$ is a configurable maximum time gap (default 30\,s).
The merged entity spans $[\min(\text{first\_seen}),
\max(\text{last\_seen})]$ and inherits the union of properties.


% ---------------------------------------------------------------------------
\subsection{Frame Embeddings and Vector Database}%
\label{sec:embeddings}

In parallel with entity extraction, the pipeline samples frames
within verified clip boundaries at a configurable rate (default
1\,fps) and extracts visual embeddings using
SigLIP-2~\cite{zhai2023siglip} (\texttt{siglip2-base-patch16-256},
768-d).
Embeddings are stored in a SQLite-backed vector database alongside
frame metadata (video ID, timestamp).

At retrieval time, a text query is encoded by the same SigLIP-2
text encoder and compared against stored frame embeddings via
cosine similarity, enabling cross-modal semantic search
(Section~\ref{sec:retrieval}).


% ---------------------------------------------------------------------------
\subsection{Agentic Retrieval}\label{sec:retrieval}

The retrieval agent enables users to query the processed video
corpus with high-level task descriptions (\eg, ``make me a
coffee'') and receive ranked clips suitable for robot policy
training.  The agent is implemented as a LangGraph
\texttt{StateGraph} supporting both sequential and parallel
execution modes (Figure~\ref{fig:retrieval-agent}).

\begin{figure*}[t]
  \centering
  \includegraphics[width=\textwidth]{retrieval_langgraph.pdf}
  \caption{%
    \textbf{Retrieval agent (parallel mode).}
    The task decomposer breaks the user query into $N$ sub-tasks
    and fans them out via the LangGraph \texttt{Send} API.
    Each sub-task runs an independent search subgraph containing
    a search planner $\leftrightarrow$ executor
    $\leftrightarrow$ analyzer loop.
    Results are aggregated by the VQA synthesiser, which selects
    and extracts the best clips for robot policy training.
  }
  \label{fig:retrieval-agent}
\end{figure*}

\subsubsection{Task Decomposition}

The user's query is first decomposed by an LLM into a set of
concrete sub-tasks, each specifying a \texttt{target\_action}
and \texttt{target\_object}.  For example:
\begin{quote}\small
  ``make me a coffee'' $\;\to\;$ \{grab mug, open coffee machine lid,
  insert capsule, close lid, press button, pour coffee\}
\end{quote}

\subsubsection{Search Planner with Strict-to-Relaxed Hierarchy}

For each sub-task, a \emph{search planner} node decides which
tool to invoke and at what relaxation level.  The entity graph
tool (\texttt{search\_graph}) supports four levels applied to
action segment queries:

\begin{enumerate}[leftmargin=*,itemsep=1pt,label=\textbf{L\arabic*}:]
  \item \textbf{Strict.} Exact match on both \texttt{action} and
    \texttt{object}.
  \item \textbf{Partial object.} Exact action, partial
    (\texttt{LIKE}) object match.
  \item \textbf{Partial both.} Partial match on action and object.
  \item \textbf{Any.} Match on action \emph{or} object, including
    free-text search in clip descriptions.
\end{enumerate}

An \emph{analyzer} node evaluates results after each search
attempt and reports back to the planner, which decides whether to
escalate to the next relaxation level, switch to the frame
embedding tool (\texttt{search\_frames}), or accept the current
results.

\subsubsection{Parallel Sub-Task Execution}

In parallel mode (default), the task decomposer emits one
LangGraph \texttt{Send} message per sub-task.  Each sub-task is
executed by an independent instance of the search subgraph
(planner $\leftrightarrow$ executor $\leftrightarrow$ analyzer
loop).  Results are aggregated back into the parent state via
reducer-annotated keys, avoiding concurrent-write conflicts.

\subsubsection{VQA Synthesiser}

The final node receives all search results, selects the best
clips for policy training (ranked by relevance and temporal
coverage), extracts the recommended video segments using FFmpeg,
and generates a structured answer.


% ---------------------------------------------------------------------------
\subsection{Scaling Architecture}\label{sec:scaling}

Processing hundreds of demonstration videos requires careful
parallelisation (Figure~\ref{fig:scaling}).
\sys{} provides the following mechanisms:

\begin{figure*}[t]
  \centering
  \includegraphics[width=\textwidth]{scaling_architecture.pdf}
  \caption{%
    \textbf{Scaling architecture.}
    Videos are discovered, probed for duration via \texttt{ffprobe},
    and distributed across GPU workers using Longest-Processing-Time
    (LPT) sharding.
    Each worker runs the full ingestion pipeline
    (segmentation, verification, entity extraction, frame embedding)
    on its shard of videos.
    All workers write concurrently to shared SQLite databases
    operating in WAL mode.
    The retrieval agent executes sub-tasks in parallel via the
    LangGraph Send API.
  }
  \label{fig:scaling}
\end{figure*}

\paragraph{Multi-GPU LPT sharding.}
Each video's duration is probed via \texttt{ffprobe} (header-only,
$<$10\,ms per video).  Videos are sorted longest-first and greedily
assigned to the GPU worker with the smallest total load---a variant
of the Longest Processing Time (LPT) scheduling algorithm.
This produces near-optimal load balance when video lengths vary
widely.

\paragraph{Chunked VLM processing.}
Rather than feeding an entire video to the VLM (which would exceed
context limits and degrade accuracy), each video is processed in
overlapping chunks.  This bounds per-call memory and latency while
the overlap ensures no action boundary is missed.

\paragraph{Resume capability.}
A progress file (\texttt{progress\_worker\_\{id\}.jsonl}) records
the outcome of each video.  On restart, already-processed videos
are filtered out via \texttt{filter\_processed\_videos}, enabling
crash-resilient batch processing.

\paragraph{Concurrent database writes.}
Multiple workers write to a single shared SQLite database by
operating in WAL (Write-Ahead Logging) mode, which permits
concurrent readers and serialised writers without lock contention
on read-heavy workloads.


% ============================================================================
% 4  BENCHMARK AND EVALUATION
% ============================================================================
\section{Benchmark and Evaluation}\label{sec:eval}

We evaluate the ingestion pipeline on
EPIC-KITCHENS-100~\cite{damen2022rescaling}, a large-scale
egocentric video dataset of kitchen activities with dense temporal
annotations.

\subsection{Evaluation Protocol}

We adopt a two-tier evaluation framework that separately measures
\emph{temporal segmentation quality} and
\emph{annotation accuracy}, reflecting the distinct failure modes
of zero-shot VLM-based segmentation.

\subsubsection{Temporal Segmentation Metrics}

\begin{itemize}[leftmargin=*,itemsep=2pt]
  \item \textbf{mAP@tIoU.}
    Mean Average Precision at temporal Intersection-over-Union
    thresholds $\{0.1, 0.2, 0.3, 0.4, 0.5\}$.
    When available, we use the official C2-Action-Detection
    \texttt{ANETdetection} evaluator for leaderboard-comparable
    scores; otherwise a standalone implementation is used.

  \item \textbf{Boundary Precision / Recall / F1.}
    For each ground-truth segment boundary (start or end), we check
    whether a predicted boundary falls within a tolerance window
    $\delta \in \{1.0, 2.0\}$\,s, and vice versa.

  \item \textbf{Segmentation ratio.}
    The ratio $|\hat{\mathcal{S}}| / |\mathcal{S}^*|$ of predicted
    to ground-truth segment counts, per video.  Values $> 1.5$
    indicate over-segmentation; $< 0.5$ under-segmentation.
\end{itemize}

\subsubsection{Annotation Accuracy Metrics}

Because the pipeline produces free-text action and object labels
rather than class indices, we employ multiple complementary metrics:

\begin{itemize}[leftmargin=*,itemsep=2pt]
  \item \textbf{Top-$k$ class accuracy.}
    Free-text predictions are mapped to the EPIC-KITCHENS verb and
    noun taxonomies via sentence-transformer cosine similarity.
    We report top-1 and top-5 accuracy for both verbs and nouns
    on temporally matched segments (tIoU $\geq 0.3$).

  \item \textbf{Soft match rate.}
    The fraction of matched segments where the ground-truth verb
    (resp.\ noun) appears as a substring of the predicted
    action (resp.\ object+description) text.

  \item \textbf{Semantic similarity.}
    Mean cosine similarity between sentence-transformer embeddings
    of predicted descriptions and ground-truth narrations
    (MiniLM-L6-v2).

  \item \textbf{BERTScore.}
    Precision, recall, and F1 of contextual token-level similarity
    between predicted descriptions and ground-truth
    narrations~\cite{zhang2020bertscore}.
\end{itemize}


\subsection{Results}

\begin{table}[t]
  \centering
  \caption{%
    \textbf{Temporal segmentation results} on EPIC-KITCHENS-100
    validation set.  Results to be populated after full benchmark
    run.
  }
  \label{tab:temporal}
  \small
  \begin{tabular}{lc}
    \toprule
    \textbf{Metric} & \textbf{Score} \\
    \midrule
    mAP@0.1 & --- \\
    mAP@0.3 & --- \\
    mAP@0.5 & --- \\
    mAP avg  & --- \\
    \midrule
    Boundary F1 @1.0\,s & --- \\
    Boundary F1 @2.0\,s & --- \\
    \midrule
    Seg.\ ratio (mean) & --- \\
    \bottomrule
  \end{tabular}
\end{table}

\begin{table}[t]
  \centering
  \caption{%
    \textbf{Annotation accuracy} on EPIC-KITCHENS-100 validation
    set (matched at tIoU $\geq 0.3$).
  }
  \label{tab:annotation}
  \small
  \begin{tabular}{lc}
    \toprule
    \textbf{Metric} & \textbf{Score} \\
    \midrule
    Verb top-1 acc  & --- \\
    Verb top-5 acc  & --- \\
    Noun top-1 acc  & --- \\
    Noun top-5 acc  & --- \\
    \midrule
    Verb soft match rate & --- \\
    Noun soft match rate & --- \\
    \midrule
    Semantic sim.\ (mean) & --- \\
    BERTScore F1          & --- \\
    \bottomrule
  \end{tabular}
\end{table}

Tables~\ref{tab:temporal} and~\ref{tab:annotation} present
placeholder entries for the temporal segmentation and annotation
accuracy results, respectively.  These will be populated after a
full benchmark run across the EPIC-KITCHENS-100 validation split.

\paragraph{Qualitative observations.}
Preliminary runs indicate that the verification--refinement loop
improves boundary precision by correcting clips where the VLM
initially over- or under-estimated action boundaries.
The LLM-based deduplication reduces false duplicates from chunk
overlaps more reliably than the heuristic baseline, at the cost
of additional LLM calls.


% ============================================================================
% 5  DISCUSSION AND FUTURE WORK
% ============================================================================
\section{Discussion and Future Work}\label{sec:discussion}

\paragraph{Verification cost vs.\ quality.}
Each refinement iteration requires additional VLM inference over
extracted clip videos.  In practice, 1--2 iterations suffice for
most videos; beyond that, diminishing returns set in.
Adaptive iteration budgets based on per-video quality estimates
could reduce cost.

\paragraph{Entity linking beyond string matching.}
The current linker merges entities by exact ID match within a
time window.  Embedding-based entity resolution (\eg, comparing
SigLIP features of object crops) would handle cases where the
LLM assigns different surface names to the same physical object
across chunks.

\paragraph{Retrieval agent grounding.}
The retrieval agent currently relies on the entity graph and
frame embeddings independently.  Fusing both modalities---\eg,
re-ranking graph results using visual similarity---could improve
retrieval precision.

\paragraph{Closed-loop policy training.}
An exciting direction is to close the loop from retrieval to
policy training: the agent selects clips, trains a policy, evaluates
it in simulation, and iteratively refines its clip selection based
on policy performance.

\paragraph{Cross-video entity resolution.}
When processing a corpus of videos filmed in the same environment,
the same physical objects appear across videos.  Cross-video entity
linking would enable richer queries (\eg, ``show all interactions
with the blue mug across all recordings'').


% ============================================================================
% 6  CONCLUSION
% ============================================================================
\section{Conclusion}\label{sec:conclusion}

We presented \sys{}, an agentic framework for automated analysis
of robot demonstration videos.  The system combines VLM-based
video segmentation with a self-correcting verification loop,
LLM-driven entity graph construction, SigLIP-2 visual embeddings,
and an agentic retrieval system with strict-to-relaxed search.
Built on LangGraph, the framework scales to large video corpora
via multi-GPU LPT sharding and concurrent database writes.
We defined a two-tier evaluation protocol on EPIC-KITCHENS-100
covering both temporal segmentation quality and annotation
accuracy, providing a rigorous framework for benchmarking
zero-shot VLM-based video understanding systems for robotics
applications.


% ============================================================================
% REFERENCES
% ============================================================================
\bibliographystyle{plain}
\begin{thebibliography}{99}

\bibitem{bai2023qwen}
J.~Bai \etal,
``Qwen-VL: A versatile vision-language model for understanding,
localization, text reading, and beyond,''
\emph{arXiv:2308.12966}, 2023.

\bibitem{brohan2023rt2}
A.~Brohan \etal,
``RT-2: Vision-language-action models transfer web knowledge to
robotic control,''
\emph{arXiv:2307.15818}, 2023.

\bibitem{chi2023diffusion}
C.~Chi \etal,
``Diffusion policy: Visuomotor policy learning via action
diffusion,''
\emph{RSS}, 2023.

\bibitem{cong2021spatial}
Y.~Cong \etal,
``Spatial-temporal transformer for dynamic scene graph
generation,''
\emph{ICCV}, 2021.

\bibitem{damen2022rescaling}
D.~Damen \etal,
``Rescaling egocentric vision: Collection, pipeline and
challenges for EPIC-KITCHENS-100,''
\emph{IJCV}, 2022.

\bibitem{ji2020action}
J.~Ji \etal,
``Action genome: Actions as compositions of spatio-temporal
scene graphs,''
\emph{CVPR}, 2020.

\bibitem{langgraph2024}
LangChain,
``LangGraph: Build resilient language agents as graphs,''
\url{https://github.com/langchain-ai/langgraph}, 2024.

\bibitem{lei2021detecting}
J.~Lei \etal,
``Detecting moments and highlights in videos via natural language
queries,''
\emph{NeurIPS}, 2021.

\bibitem{li2023videochat}
K.~Li \etal,
``VideoChat: Chat-centric video understanding,''
\emph{arXiv:2305.06355}, 2023.

\bibitem{lin2023video}
B.~Lin \etal,
``Video-LLaVA: Learning united visual representation by
alignment before projection,''
\emph{arXiv:2311.10122}, 2023.

\bibitem{nvidia2025cosmos}
NVIDIA,
``Cosmos: World foundation models for physical AI,''
\emph{Technical Report}, 2025.

\bibitem{schick2023toolformer}
T.~Schick \etal,
``Toolformer: Language models can teach themselves to use
tools,''
\emph{NeurIPS}, 2023.

\bibitem{shi2023tridet}
D.~Shi \etal,
``TriDet: Temporal action detection with relative boundary
modeling,''
\emph{CVPR}, 2023.

\bibitem{wu2023autogen}
Q.~Wu \etal,
``AutoGen: Enabling next-gen LLM applications via multi-agent
conversation,''
\emph{arXiv:2308.08155}, 2023.

\bibitem{xu2017video}
D.~Xu \etal,
``Video question answering via gradually refined attention over
appearance and motion,''
\emph{ACM MM}, 2017.

\bibitem{xu2025egagent}
W.~Xu \etal,
``EGAgent: An adaptive ego-centric video understanding agent,''
\emph{arXiv:2601.18157}, 2025.

\bibitem{yao2023react}
S.~Yao \etal,
``ReAct: Synergizing reasoning and acting in language models,''
\emph{ICLR}, 2023.

\bibitem{zhai2023siglip}
X.~Zhai \etal,
``Sigmoid loss for language image pre-training,''
\emph{ICCV}, 2023.

\bibitem{zhang2020bertscore}
T.~Zhang \etal,
``BERTScore: Evaluating text generation with BERT,''
\emph{ICLR}, 2020.

\bibitem{zhang2022actionformer}
C.~Zhang \etal,
``ActionFormer: Localizing moments of actions with
transformers,''
\emph{ECCV}, 2022.

\bibitem{zhao2023aloha}
T.~Z.~Zhao \etal,
``Learning fine-grained bimanual manipulation with
low-cost hardware,''
\emph{RSS}, 2023.

\end{thebibliography}

\end{document}
